\chapter{分块与莫队}

\section{序列分块基础}

序列分块的基本想法是把序列按照某个大小 $B$ 分段（块）。任何一个区间 $[l,r]$ 在这个结构上要么是某个块的一个子段，要么是一个块的后缀加 $O(n/B)$ 个完整块加一个块的前缀。这种结构的好处在于，我们通过\textbf{重构}而不是合并来维护信息：在线段树中，当某个结点的子树内部发生了某些修改时，必须要能够通过某种方式正确地计算出此结点上新的信息；而在分块结构中，我们只要根据这个块对应的 $O(B)$ 长度的序列重新计算这个结点的信息就可以了，也就是\textbf{重构}。当然，对于整块的修改，还是需要能够快速地把修改标记应用到信息上。

根据是否存在修改可以将问题分为动态问题与静态问题。对于动态问题，我们通常要按顺序思考以下几个事情：
\begin{itemize}
    \item 把块长设为\textbf{待定值} $B$。
    \item 如何查询一个完整块的答案？通常来说，这通过对每个完整块维护一个数据结构来实现。
    \item 如何查询一个零散块的答案？通常使用暴力，或者聪明的暴力。
    \item 如何在完整块修改的时候快速地更新完整块的数据结构？通常通过打某种标记实现。此外，也许需要改变第一条中的使用的数据结构，以保证在修改时可以快速维护。
    \item 如何在零散块修改的时候更新对应完整块的数据结构？通常通过\textbf{暴力重构}实现。
    \item 最后，计算你的复杂度，并选择最优的 $B$。
\end{itemize}

在实践中，块长 $B$ 也可以通过实验的方式得出，并且这样得到的实际效率也许更高。相比分治数据结构，分块更加灵活，需要思考的事情也更多。我们先前和以后学到的各种思想和技巧也经常在分块中使用。

\begin{example}
    区间加，查询区间内小于 $x$ 的数的个数。
\end{example}
\begin{remark}
    思考一下为什么线段树做不了这个。事实上，即使是区间 $\pm 1$，查询全局 $0$ 的个数也是无法 $\mathrm{polylog} n$ 的。
\end{remark}
\begin{solution}
    \begin{itemize}
        \item 查询完整块的答案：维护排序后的数组，查询时二分。$O(\log B)$ 查询一个块，总共有 $qn/B$ 次查询。
        \item 查询零散块的答案：遍历即可。$O(B)$ 查询一个块，总共有 $q$ 次查询。
        \item 完整块的修改：打一个加法标记，$O(1)$。容易发现这与第一条中的数据结构是兼容的。总共有 $qn/B$ 次修改。
        \item 零散块的修改：重构。直接排序是 $O(B\log B)$ 的，但我们可以注意到可以复用先前排序的数组。把修改过和没修改过的数分开，归并这两类数即可做到 $O(B)$。总共有 $q$ 次修改。
        \item 计算复杂度：$(qn/B)\log B+qB$。可以这样计算复杂度：首先，可以猜想 $B=n^{\Theta(1)}$，所以 $\log B=\Theta(\log n)$，这样复杂度就变成 $(qn/B)\log n+qB$。接下来，由于两个组成部分的单调性不同，其最优值在渐近复杂度的意义下必然在两部分相等时取到，从而只需要解方程 $(qn/B)\log n=qB$，即 $B=\sqrt{n\log n}$。这样得到复杂度 $q\sqrt{n\log n}$.
    \end{itemize} 
\end{solution}
\begin{remark}
    有比较复杂的严格根号做法，但是不太实用。
\end{remark}


\begin{example}[P5356]
    区间加，查询区间第 $k$ 小。
\end{example}
\begin{solution}
    首先可以在上一题的基础上套一层二分得到 $O(\sqrt{n\log n}\log V)$ 的做法.

    注意到二分的时候并不需要每次都花费 $B$ 的代价查询散块，而是可以先将散块归并成一个临时的整块，查询时在上面二分。这样散块的代价就从 $B\log V$ 降到 $B+\log B\log V$，总复杂度 $O(\sqrt{n\log n\log V})$。
\end{solution}

\begin{exercise}[P3863]
    区间加，查询某一个位置在过去的多少个时间内 $\le y$（$y$ 每次不同）
\end{exercise}

\begin{example}[行加列和]
    给定 $n$ 个点，支持把 $x_i\in [X_l,X_r]$ 的点的权值加上一个数，查询 $y_i\in [Y_l,Y_r]$ 的点权和。
\end{example}
\begin{solution}
    不失一般性，可以假设所有的点的 $x,y$ 都不同，并离散化询问与修改区间。按 $x$ 分块。
    \begin{itemize}
        \item 查询：枚举每个块，查询此块内 $y$ 在 $[Y_l,Y_r]$ 的点的权值和。这可以差分后用前缀和 $s_{i,j}$ 完成：$s_{i,j}$ 表示第 $i$ 个块内 $y$ 排名前 $j$ 的点的权值和。为了 $O(1)$ 得到 $j$，还需要预处理一个数组 $\mathrm{rank}_{i,y}$ 表示第 $i$ 个块内 $y_j\le y$ 的点的个数。
        \item 完整块修改：打一个加法标记。在询问时，还需要额外加上标记值乘 $\mathrm{rank}_{i,y}$。
        \item 散块修改：重算 $s_{i}$ 即可。
    \end{itemize}
    这个做法需要 $O(n\sqrt{n})$ 的空间。
\end{solution}
\begin{solution}[压空间]
    由于不同块的修改和查询独立，可以离线后逐块处理，这样只需要存储一个块的 $\mathrm{rank}_y$，空间降至线性。

    另一种方式是对 $y$ 也分块，查询时单独查询散块的答案。
\end{solution}
\begin{solution}[另解]
    是上一个做法的对偶。按行列分别分块。差分掉修改。
    \begin{itemize}
        \item 查询完整 $y$ 块的答案：维护 $y$ 在此块内的点的权值和
        \item 查询零散 $y$ 块的答案：暴力。
        \item 前缀的若干个完整 $x$ 块的修改：首先对每个 $x$ 块打一个标记表示此块上有一个修改（用于零散 $y$ 块的查询），然后考虑如何更新 $y$ 块维护的答案。对每个 $y$ 块，预处理一个系数 $C_{y,x}$，表示当前 $x$ 个 $x$ 块发生了加一的时候，此 $y$ 块的答案会变化多少。据此更新所有 $y$ 块的答案即可。
        \item 零散 $x$ 块的修改：暴力修改点权。
    \end{itemize}
\end{solution}

\begin{exercise}[CodeChef Chef and Churu]
    维护一个序列 $a$ 和一个函数列 $f$，每个函数都是 $a$ 的一个区间和。支持：序列的单点修改；查询函数列的区间和。
\end{exercise}

\begin{example}[区间众数]
    查询区间内出现次数最多的数的出现次数。
\end{example}
\begin{solution}
    首先用 $O(n^2/B)$ 的时间 计算出数组 $M_{x,y}$ 表示第 $x$ 到第 $y$ 个块中出现次数最多的数的出现次数。查询时，还需要枚举散块中的数判断能否更新答案。具体来说，每次需要做 $O(B)$ 次判断：值 $w$ 在 $[l,r]$ 中的出现次数是否至少是 $t$。聪明的办法是：对每个值维护一个 \texttt{vector} 存储其所有出现，则只需要找到目前枚举的 $w$ 在这个 \texttt{vector} 中出现的位置 $p$，然后判断 \texttt{vector} 的 $p+t$ 位置是否 $\le r$ 即可。
\end{solution}

\begin{example}[区间逆序对]
    查询区间内逆序对的数量。
\end{example}
\begin{solution}
    首先用 $O(n^2/B)$ 的时间 计算出数组 $M_{x,y}$ 表示第 $x$ 到第 $y$ 个块中的逆序对数量。查询时，还需要加上散块的贡献。聪明的办法是：首先处理一个散块内部的贡献，这可以预处理；然后处理两个散块之间的贡献，这可以通过归并实现；最后处理散块与完整块之间的贡献，这相当于 $B$ 次查询：第 $x$ 到第 $y$ 个块中有多少个数 $<w$，可以差分后预处理。

    但是这还遗漏了一种情况：询问区间完全被包含在某个块中。这个情况可以通过差分变成两个散块前缀之间的询问。
\end{solution}
\begin{remark}
    未知是否存在 $\tilde O(n)$ 空间的做法，但是已知使用 $O(1)$ 二维数点来回答散块询问的话一定做不到 $\tilde O(n)$ 空间。
\end{remark}
\begin{remark}
    以上两个问题可以视为在线莫队。
\end{remark}


\section{复杂度平衡}

线段树可以实现 $O(\log n)$ 的单点修改和区间查询。然而，有时修改和查询的数量是不同的。此时，可以使用多叉树来降低某一侧的复杂度而提高另一侧的复杂度，从而平衡整体的复杂度。

\begin{example}
    使用 $B$ 叉树分别实现：
    \begin{itemize}
        \item $O(B\log_B n)$ 单点加，$O(\log_B n)$ 查询区间和。
        \item $O(\log_B n)$ 单点加，$O(B\log_B n)$ 查询区间和。
    \end{itemize}
\end{example}
\begin{exercise}[对偶]
    实现
    \begin{itemize}
        \item $O(B\log_B n)$ 区间加，$O(\log_B n)$ 查单点。
        \item $O(\log_B n)$ 区间加，$O(B\log_B n)$ 查单点。
    \end{itemize}
\end{exercise}
\begin{exercise}[差分]
    实现
    \begin{itemize}
        \item $O(B\log_B n)$ 区间加，$O(\log_B n)$ 查询区间和。
        \item $O(\log_B n)$ 区间加，$O(B\log_B n)$ 查询区间和。
    \end{itemize}
\end{exercise}
\begin{exercise}
    分别代入 $B=\log n$，$B=\sqrt{n}$ 和 $B=n^{\epsilon}$，观察以上几个问题的复杂度。
\end{exercise}
\begin{exercise}
    分别实现数据结构，支持
    \begin{itemize}
        \item $O(V^{\epsilon}/\epsilon)$ 插入删除，$O(1/\epsilon)$ 查询 $<x$ 的数的数量（即 $x$ 的排名）。
        \item $O(1/\epsilon)$ 插入删除，$O(V^{\epsilon}/\epsilon)$ 查询 $<x$ 的数的数量。
    \end{itemize}
\end{exercise}


\begin{example}
    使用 $B$ 叉树实现：$O(B^2\log_B n)$ 单点修改，$O(\log_B n)$ 查询区间分治信息和（假设分治信息可以 $O(1)$ 合并）。
    \begin{hint}
        $B$ 叉树可以将问题变成 $\log_B n$ 个大小为 $B$ 的子问题。
    \end{hint}
\end{example}
\begin{remark}
    在修改的时候，还需要能够重新计算每个结点的 $B$ 个孩子的总信息，这往往使得修改难以优于 $B\log_B n$。\href{https://www.cnblogs.com/-Houraisan-Kaguya/p/19601799}{这篇文章}证明了对于一般的半群信息，在常规的模型下，如果修改复杂度是 $O(f(n))$，那么无法做到 $\frac{n}{e^{O(f)}}$ 的全局查询。
    
    不过，在信息比较特殊的情况下，还是有可能利用位运算之类的 word RAM 特性加速的，例如 $O(\log n/\log w)$ 的单点加一查询区间和，以及先前见到的线性 RMQ。此外，通过非常复杂的方法还可以做到 $O(\log n/\log\log n)$ 的单点修改 RMQ。
\end{remark}

\begin{example}
    维护一个多重集，分别支持
    \begin{itemize}
        \item $O(\log_B V)$ 插入删除，$O(B\log_B V)$ 查询最大值。
        \item $O(B\log_B V)$ 插入删除，$O(\log_B V)$ 查询最大值。
    \end{itemize}
\end{example}

\begin{example}
    维护一个\textbf{多重集}，分别支持
    \begin{itemize}
        \item $O(\log_B V)$ 插入删除，$O(B\log_B V)$ 查询第 $k$ 小。
        \item $O(B\log_B V)$ 插入删除，$O(\log_B V)$ 查询第 $k$ 小。
    \end{itemize}
\end{example}
\begin{solution}
    对于前者，在 $B$ 叉树上类似二分地做即可。

    对于后者，在 $B$ 叉树的每个结点上维护所有第 $k$ 小所在的孩子。修改时只需要更改 $B$ 个位置。
\end{solution}

\begin{example}
    维护一个序列，分别支持
    \begin{itemize}
        \item $O(B\log_B n)$ 单点修改，$O(\log_B n)$ 查询区间最值。
        \item $O(\log_{B_1} n \log_{B_2}V)$ 单点修改，$O((B_1\log_{B_1} n)(B_2\log_{B_2} V))$ 查询区间最值。（这没有那么简单）
    \end{itemize}
\end{example}
\begin{solution}
    值域树套区间树。
\end{solution}


\section{莫队}

莫队算法可以视为一般的扫描线：我们用数据结构维护当前的某个集合 $S$，通过加删元素来在不同的 $S$ 之间转移。形式化的描述如下：

\begin{definition}[莫队]
    给定一个集合族 $S_1,\cdots,S_m$。你有一个初始为空的集合 $S$，每次可以进行以下操作之一：
    \begin{itemize}
        \item 向 $S$ 中添加一个元素
        \item 删除 $S$ 中的某个元素
        \item 报告 $i$ 表示当前的 $S$ 是 $S_i$。
        \item 结束操作
    \end{itemize}
    需要保证在结束操作前，每个 $S_i$ 都被正确报告过。莫队算法的目标是构造一个尽量短的操作序列。
\end{definition}

莫队问题有一些常见的变体：

\begin{definition}[不删除莫队/回滚莫队]
    在莫队问题中，不允许删除元素，而是改为允许撤销上一次添加的元素。
\end{definition}
\begin{definition}[在线莫队]
    $S_i$ 变成在线给定的，而操作额外允许记录某一个版本以及转移到某个记录的版本上。
\end{definition}

\begin{example}[前缀莫队/扫描线]
    每个 $S_i$ 是一个前缀 $[1,i]$。
\end{example}
\begin{remark}
    这自然是不删除莫队。
\end{remark}

接下来，我们给出常见的几种莫队构造。

\begin{example}[区间莫队]
    每个 $S_i$ 是一个区间 $[l_i,r_i]$。
\end{example}
\begin{solution}
    将序列按照 $B$ 分块，按照 $(l_i\text{ 所在块},r_i)$ 排序。对于 $l_i$ 在同一个块内的询问，转移代价不超过 $B$；$l_i$ 所在块至多变化 $n/B$ 次，每次变化的代价至多为 $n$，所以总代价不超过 $qB+n^2/B$。取 $B=n/\sqrt{q}$ 时得到最优复杂度 $O(n\sqrt{q})$。
\end{solution}
\begin{remark}
    注意这里的块长并不是常规的 $\sqrt{n}$。事实上，如果取 $B=\sqrt{n}$ 的话可以得到 $(n+q)\sqrt{n}$ 的复杂度，这比 $n\sqrt{q}$ 严格更差。
\end{remark}
\begin{exercise}
    给出在线不删除莫队的构造。
\end{exercise}

\begin{example}[单点补莫队]
    每个 $S_i$ 是元素 $i$ 的补集。
\end{example}
\begin{exercise}
    说明区间补莫队与区间莫队等难。
\end{exercise}

\begin{example}[单点补不删除莫队/线段树分治]
    同上，但是需要不删除莫队。
\end{example}
\begin{solution}
    DFS 线段树，进入一个结点的时候将其兄弟结点对应的区间加入，退出结点的时候撤销。
\end{solution}
\begin{exercise}
   假设 $K=O(1)$，构造所有长度不超过 $K$ 的区间的补集的不删除莫队。     
\end{exercise}


\begin{example}[对偶区间莫队/线段树分治]
    每个元素是一个区间，而 $S_i$ 是所有包含 $i$ 的区间的集合。
\end{example}
\begin{solution}
    DFS 线段树，进入一个结点后保证所有覆盖此结点对应区间的区间都在集合中。则把所有区间在线段树上分解，DFS 到一个结点时加入此结点上的区间即可。

    可以一边 DFS 一边分解区间来将空间降到线性。
\end{solution}
\begin{exercise}
    对于允许离线的问题，可以通过这种方式将删除变为撤销。一个经典的例子是\textbf{动态图连通性}，即加删边，查询两点连通性。
\end{exercise}
\begin{exercise}
    证明不删除莫队的代价至多比普通莫队多一个 $\log$。
\end{exercise}


\begin{example}[子树莫队/dsu on tree]
    给定一棵树，每个 $S_i$ 是一个子树。
\end{example}
\begin{exercise}
    函数 \texttt{solve(u,0/1)} 表示构造出 $u$ 子树内所有子树对应的操作序列，且函数结束后是否保留此子树内点，$0/1$ 分别对应不保留/保留。首先树链剖分，然后调用 \texttt{solve(root,1)}。$\mathtt{solve(u,0/1)}$ 具体实现如下：首先递归所有轻儿子 $\mathtt{solve(v_l,0)}$，然后递归重儿子 $\mathtt{solve(v_h,1)}$，最后加入 $u$ 及其所有轻子树。如果不保留 $u$ 子树内的点，则撤销 $\mathtt{size_u}$ 次即可。总复杂度为所有点的轻子树大小之和，这是 $O(n\log n)$ 的。
\end{exercise}

\begin{theorem}[莫队之理]
    设 $P'(i)$ 为 $i$ 个点能够划分出的最多询问等价类数量，则莫队的最优复杂度是 $\tilde \Theta(n\max_{i}P'(i)/i)$。这个最优复杂度通过可以如下实现近似：找到 $x$ 使得 $P'(x)\approx m$，随机选取 $x$ 个点，对每个集合 $S_i$ 构造 01 向量 $B_i$，其中 $B_{i}(j)=[j\in S_i]$。按照这些向量的字典序排序即可得到一个合适的顺序，此顺序可以保证总代价高概率不超过最优代价乘 $\log n\log m$。
\end{theorem}
\begin{proof}
    道理比较复杂，在此省略。可以参考 \href{https://www.luogu.com.cn/article/o6pe9wm3}{critnos 的博客}。
    % 考虑这 $x$ 个点划分出的至多 $m$ 个询问等价类。考虑每个点对操作次数造成的贡献。对于排序后相邻的询问，只有 $x$ 恰好在其中一个询问集合中出现时，才会造成一次贡献。
    
    % 对每个等价类，
\end{proof}
\begin{theorem}
    对于在线莫队，只需在每个等价类中任取一个集合加入需要预处理的集合即可。在同一个等价类内，无论如何转移，复杂度都是正确的。
\end{theorem}

在完成了莫队操作序列的构造之后，我们就可以使用数据结构来维护这些操作了。对于静态问题，只要不强制在线，莫队就总是好于分块。

\begin{example}[小 Z 的袜子]
    查询区间内 $a_i=a_j$ 的点对数量。另一种表述是每种数出现次数的平方和。
\end{example}
\begin{remark}
    存在基于矩阵乘法的 $O(n^{\frac{2\omega}{\omega+1}})$ 做法
\end{remark}

\begin{example}[区间众数]
    查询区间内出现次数最多的数的出现次数。
\end{example}
\begin{remark}
    存在 $O(n^{1.49})$ 做法。
\end{remark}

\begin{example}[T142605]
    查询区间内出现次数第 $k_1$ 小的数里面的第 $k_2$ 小。
\end{example}

\begin{example}[区间逆序对]
    查询区间内逆序对的数量。
\end{example}
\begin{solution}
    使用区间莫队，转移时需要查询当前集合内有多少个 $a_i<x$。直接用树状数组做是 $O(n\sqrt{m}\log n)$ 的。

    注意贡献可以差分，即每次都询问 $[1,r]$ 内有多少个 $a_i<x$。这样问题变成在线的 $O(n\sqrt{m})$ 次二维数点，可以通过可持久化的 $B$ 叉树做到 $O(1)$ 询问。

    注意到并不需要真的在线。每一批贡献都形如
    $$
    \sum_{i=r+1}^{r'}([l,i-1]\text{ 内满足 }a_j<a_i\text{ 的数的数量}).
    $$
    差分后，就变成
    $$
    \sum_{i=r+1}^{r'}([1,l-1]\text{ 内满足 }a_j<a_i\text{ 的数的数量})
    $$
    和
    $$
    \sum_{i=r+1}^{r'}([1,i-1]\text{ 内满足 }a_j<a_i\text{ 的数的数量}).
    $$
    对于前者，可以用一个 $(r,r',l-1)$ 的元组来表示一批贡献；对于后者，可以用一个 $(r,r')$ 的元组来表示一批贡献。对于这 $O(m)$ 个元组重新离线即可去掉可持久化。
\end{solution}
\begin{remark}
    存在在线的基于矩阵乘法的 $O(n^{1.41})$ 做法。
\end{remark}

\begin{example}[P4869]
    给一棵树，记 $S_u$ 为 $u$ 的子树，查询 $\sum_{i\in S_u}\sum_{j\in S_v}[a_i=a_j]$。
\end{example}
\begin{remark}
    这是伪双子树莫队。真双子树莫队是 P9988。
\end{remark}

\begin{example}[查询区间最小差]
    查询区间内最小的 $|a_i-a_j|$（$i\ne j$）
\end{example}
\begin{remark}
    这个其实可以 polylog，参见 CF765F。
\end{remark}

\begin{example}[k 区间颜色数]
    给一个序列，每次给出k个区间，求k个区间的并内有多少个种出现过的值。$O(nm/w)$
\end{example}

\section{常用技巧}

这里是一些常用的技巧，但是暂时没有合适的例题来展示。

\subsection{逐块处理}

如果可以离线，并且块之间的贡献相对独立，那么将所有修改和查询离线后\textbf{逐块处理}总是更好的选择。

\subsection{操作分块/定期重构}

如果操作变多时处理的代价变大（如 $B^2$），但是重构的代价不随操作数量增加，那么可以考虑每处理一批操作就进行一次重构。

\subsection{分块后分治}

与上一条类似，如果一个长为 $n$ 的序列上直接做会带来较高的代价（如 $n^2$），但是可以把序列分成对询问相对独立的块，那么可以考虑先对序列做一次分块。

\begin{example}[P5611]
\end{example}